\section{Strengths and Weaknesses}
\begin{large}
\textbf{Strengths:}
\end{large}

\begin{itemize}
    \item \textbf {Data processing and sreening.} Although it is difficult to analyze text and images, we still succeed in effectively filtering reports based on text and images.
    \item \textbf {Divergent thinking.} A powerful machine learning system relies on a large number of correct samples, which we believe is not suitable for this problem. So we divergently rely on the help of the existing image recognition system, and minimize its errors from multiple dimensions.
    \item \textbf {Visualization.} We vividly and diversely display various results and ideas for thinking about problems. And, we use visualization to show the source of some data extraction, such as in the AHP model.
    \item \textbf {Robust model.} Our PS-CA model achieves a high degree of fit simulation result with only a small number of correct samples.
    \item \textbf {Solve from a multidisciplinary perspective.} Based on the problem background, we use the knowledge of population growth in biology to formulate rules for the model.
\end{itemize}

\begin{large}
\textbf{Weaknesses:}
\end{large}

\begin{itemize}
    \item In our model, the calculation of some synthesized variables is subjective, lacks credibility, and requires a lot of data to verify.
    \item Our model does not consider that pests will spread over longer distances due to human activities\parencite{simulation}.
    \item Our spread model only considers the vegetation in the area when determining the environmental factors that affect the spread of pest. If there is sufficient time and advanced technology, we can also perform niche modeling based on climate and precipitation.
    \item It is a pity that we are unable to rank all 4441 reports due to time constraints, so we finally failed to get a statistically correct rate.
\end{itemize}